% ********** Приклад оформлення пояснювальної записки **********
% *********  до атестаційної роботи ступеня бакалавра **********


\documentclass{bachelor_thesis}

% Додаткові пакети вносіть у цей файл
\input{preamble/01_packages}
\addbibresource{thesis.bib}

% Додаткові визначення та перевизначення команд вносіть у цей файл
\input{preamble/02_redefinitions}

% Бюрократичні відомості про автора роботи
%%% Основні відомості %%%
\newcommand{\UDC}                      % УДК
{(впишіть правильний УДК!)}            % УДК виглядає приблизно як 004.056.5 або 513.2, або навіть 004.056.5:513.2+519.1
% Для того, щоб знайти правильний УДК, використовуйте каталог https://teacode.com/online/udc/

\newcommand{\reportAuthor}             % ПІБ автора повністю
{Жуковін Олександр Віталійович}
\newcommand{\reportAuthorShort}        % ПІБ автора коротко
{Олександр ЖУКОВІН}
\newcommand{\reportAuthorGroup}        % група автора
{ФІ-21}
\newcommand{\reportTitle}              % Назва роботи
{Прогнозування інфляції України за допомогою методів машинного навчання}
%% використовуйте символ "\par" або "\\" для розбиття назви на декілька рядків

\newcommand{\supervisorFio}            % Науковий керівник, ПІБ повністю
{Чернятевич Антон Андрійович}
\newcommand{\supervisorFioShort}       % Науковий керівник, ПІБ коротко
{Антон ЧЕРНЯТЕВИЧ}
\newcommand{\supervisorRegalia}        % Науковий керівник: посада, степінь, звання
{доцент кафедри ММАД}              % наприклад: доцент кафедри ПЕКЛА, д.ф.-м.н., доцент
                                       % якщо виходить дуже довго - скорочуйте: доц. каф. ПЕКЛА, д.ф.-м.н., доц.

\newcommand{\consultFio}               % Консультант, ПІБ повністю
{}
\newcommand{\consultRegalia}           % Консультант: звання, степінь, посада
{}
% Якщо у вас нема консультанта - залишайте ці поля порожніми


\newcommand{\reviewerFio}              % Рецензент, ПІБ повністю
{Прізвище Ім'я По-батькові}
\newcommand{\reviewerRegalia}          % Рецензент: звання, степінь, посада
{посада, степінь, звання}

\newcommand{\YearOfDefence}            % рік захисту
{2026}
\newcommand{\YearOfBeginning}          % попередній рік - може, можна це якось автоматизувати, нє?
{2025}



% Починаємо верстку документа
\begin{document}

\pagestyle{plain}
\setfontsize{14}

% Створюємо титульну сторінку
\input{frontmatter/a00_title}

% Створюємо завдання
\input{frontmatter/a01_specification}

% Створюємо анотації
%\setcounter{page}{4}
\input{chapters/a1_abstract}

% Створюємо зміст
%\pagenumbering{gobble}
\tableofcontents
\cleardoublepage
%\pagenumbering{arabic}
%\setcounter{page}{8}    %!!! -- продумати, як автоматизувати номер сторінки

% Створюємо перелік умовних позначень, скорочень і термінів
% Якщо цей розділ вам не потрібен, просто закоментуйте два наступних рядка
\shortings
\input{chapters/a2_shortings}

% Створюємо вступ
\intro
\input{chapters/a3_introduction}

% Додаємо глави
% Якщо ваша робота містить менше або більше глав - модифікуйте наступні
% рядки відповідним чином
\input{chapters/c1_chapter_01}
\input{chapters/c2_chapter_02}
\input{chapters/c3_chapter_03}


% Створюємо висновки
\conclusions
\input{chapters/w1_conclusions}

% Додаємо бібліографію
%%%%%%% Якщо ви не використовуєте bibtex, оце все треба закоментувати звідси...
\defbibenvironment{bibliography}
  {\list
     {\printfield[labelnumberwidth]{labelnumber}}
     {\setfontsize{14}%
     }}
  {\endlist}
  {\item}

\printbibliography[heading=bibintoc,title={Перелік посилань}]
%%%%%%% ...по осюди

% Якщо у вас виникли проблеми з бібтехом (у мене, наприклад, виникли), то замість бібтеху
% можна використати менш приємний, але завжди працюючий спосіб: вписати та оформити бібліографію вручну
% Для цього закоментуйте команду \printbibliography та розкоментуйте наступний рядок
% VVV оцей рядок розкоментуйте VVV
%\input{chapters/w2_bibliography}
% ^^^ оцей-оцей ^^^
% Бібліографія вноситься у файл "w2_bibliography.tex"

%%% А це взагалі перестало працювати :(
%\bibliographystyle{ugost2008}
%\bibliography{thesis}


% Створюємо додатки (дивись у файли додатків для необхідних пояснень)
% Якщо ви маєте меншу або більшу кількість додатків, модифікуйте наступні
% рядки відповідним чином
% Якщо ви не маєте додатків, просто закоментуйте наступні рядки
\input{chapters/z1_appendix_A}
\input{chapters/z2_appendix_B}


% Нарешті
\end{document}
